\documentclass{ximera}
\title{Quiz Problems and Solutions F26}

\newcommand{\pskip}{\vskip 0.1 in}

\begin{document}
\begin{abstract}
Summary of each class.
\end{abstract}
\maketitle

\section{Quiz 1A}

\begin{question} \label{Qmfdew9e}
Find an equation of the circle with center $(-4,1)$ through the point $(-1,-3)$. 
\begin{itemize}

\item Include an accurate picture on a coordinate system. 


\item Include a complete explanation in full sentences of the \emph{logic} behind your equation. Use the Pythagorean theorem in your explanation, not the distance formula.


Note: Explaining the logic is not the same as explaining how to use a formula. It requires a lot more.

\end{itemize}

\end{question}

\section{Quiz1B}

\begin{question} \label{Qpfo3r}
Find an equation of the smallest circle through the points with coordinates $(3,-1)$ and $(-7,-9)$. 
\begin{itemize}

\item Include an accurate picture on a coordinate system. Label the axes with the appropriate variable names. Put scales on the axes.


\item Include a complete explanation in full sentences of each step in your solution. Do not save the explanation to the end.

\item Explain the \emph{logic} behind your equation. Use the Pythagorean theorem in your explanation, not the distance formula.

Note: Explaining the logic is not the same as explaining how to use a formula. It requires a lot more.

\end{itemize}

\begin{explanation}
I'll leave the pictures to you, but the key points are as follows:

\begin{enumerate}
\item Since the longest chord of a circle is a diameter, the smallest circle through the points $A(3,-1)$ and $B(-7,-9)$ is the circle that has segment $\overline{AB}$ as a diameter.

\item The center of that circle is the midpoint $M$ of $\overline{AB}$. To get the coordinates of $M$ we average the coordinates of $A$ and $B$. The midpoint therefore has coordinates
\[
  (x,y) = \left( \frac{1}{2}\left(3+ -7  \right) , \frac{1}{2}\left(-1+ -9  \right) \right) = (-2.,-5).
\]

\item The radius of the smallest circle is the distance between $M$ and either of the points $A$ or $B$. We'll find the distance between $M$ and $A$ by using the Pythagorean theorem. To do this we draw a horizontal line through $M$ and a vertical line through $A$. These lines intersect in the point $C$ with coordinates $(3,-5)$. Then the legs of right triangle $ACM$ have lengths
\[
    a = \text{dist}(M,C) = |3 - (-2)| = 5
\]
and
\[
   m = \text{dist}(A,C) = |-5 - (-1) | = 4 .
\]
So by the Pythagorean theorem the hypotenuse $\overline{AM}$ has length
\[
    c =  \sqrt{a^2 + m^2} = \sqrt{25 + 16} = \sqrt{41}.
\]

\item We can now write an equation of the circle centered at $M$ with radius $\sqrt{41}$. The key idea behind the equation is this.
 
\pskip

A point $P$ lies on the circle if and only if the distance from $P$ to the center $M$ equals the radius $\sqrt{41}$. 

\pskip

To turn this desription of the circle into an equation, we let $P$ have coordinates $(x,y)$. Then by drawing a right triangle and using the Pythagorean theorem as we did in part (c), the distance from $P(x,y)$ to the center $M(-2,-5)$ is
\[
      \sqrt{(x+2)^2 + (y+5)^2}.
\]

So a point $P$ with coordinates $(x,y)$ lies on the circle of radius $\sqrt{41}$ centered at $M(-2,-5)$ if and only if
\[
  \sqrt{(x+2)^2 + (y+5)^2} = \sqrt{41}.
\]

This is an equation of the smallest circle through the points $A(3,-1)$ and $B(-7,-9)$.

\end{enumerate}

\emph{Comment:} The most common error was failing to include an explanation of the \emph{logic} behind the equation of the circle as in part (d).
\end{explanation}
\end{question}



\end{document}